\documentclass[reqno]{amsart}
\usepackage[T1]{fontenc}
\usepackage{amsmath,amssymb,amsthm}
\usepackage{xr-hyper}
\usepackage[hidelinks]{hyperref}
\externaldocument[da-]{20261126T030300--the-data__workbook}
\externaldocument[ob-]{20261126T040100--spaces-forms-and-twists__workbook}
\externaldocument[dl-]{20261126T030500--divisibility-of-the-composite-lift__workbook}
\externaldocument[ss-]{20261126T030400--the-seesaw-identity__workbook}
\externaldocument[d2-]{20261126T050200--one-lift-at-a-small-non-split-prime__workbook}
\externaldocument[sd-]{20261125T000000--slope-method__workbook}

% Chapter number; \ref{chap} gives it.
\newcounter{chap}
% Chapter numbers, generated from 20261125T012000--grh-workbook__guide.org
% by scripts/make-chapters.py.  Do not edit: change the order of the
% index and rerun the script (./preview-site.sh and CI do this).
%
% Usage in a sheet, after \newcounter{chap}:
%   % Chapter numbers, generated from 20261125T012000--grh-workbook__guide.org
% by scripts/make-chapters.py.  Do not edit: change the order of the
% index and rerun the script (./preview-site.sh and CI do this).
%
% Usage in a sheet, after \newcounter{chap}:
%   % Chapter numbers, generated from 20261125T012000--grh-workbook__guide.org
% by scripts/make-chapters.py.  Do not edit: change the order of the
% index and rerun the script (./preview-site.sh and CI do this).
%
% Usage in a sheet, after \newcounter{chap}:
%   \input{chapters}\chapnumber{<sheet ID>}
\makeatletter
\providecommand\chapters@set[2]{\@namedef{chapters@#1}{#2}}
\chapters@set{20261125T001000}{1}% zeros-to-ramp-sums
\chapters@set{20261125T001100}{2}% inert-primes-and-estimate-b
\chapters@set{20261125T005000}{3}% galois-fourier-and-averaging
\chapters@set{20261125T005500}{4}% columns-from-a-sphere
\chapters@set{20261125T000000}{5}% slope-method
\chapters@set{20261126T030100}{6}% a-twist-for-each-prime
\chapters@set{20261125T003000}{7}% the-chain-and-the-square-root
\chapters@set{20261125T010000}{8}% rank-lemma-and-zero-estimate
\chapters@set{20261125T010500}{9}% poincare-series-and-near-orthogonality
\chapters@set{20261124T223000}{10}% elliptic-curves-for-the-proof
\chapters@set{20261125T009500}{11}% fock-model-and-theta-intertwiners
\chapters@set{20261124T223100}{12}% complex-multiplication-and-inert-reduction
\chapters@set{20261125T008000}{13}% one-cm-determinant
\chapters@set{20261124T214500}{14}% supersingular-theta-divisibility
\chapters@set{20261124T214600}{15}% the-isotropic-plane
\chapters@set{20261125T013700}{16}% pairings-on-a-formal-group-of-height-two
\chapters@set{20261125T013500}{17}% hecke-nondegeneracy-on-the-tree
\chapters@set{20261126T030200}{18}% the-bound-is-attained
\chapters@set{20261125T009000}{19}% kronecker-limit-and-chowla-selberg
\chapters@set{20261125T002000}{20}% definite-theta-lifts
\chapters@set{20261125T015500}{21}% siegel-weil-in-the-smallest-case
\chapters@set{20261125T002100}{22}% definite-theta-lifts-and-norms
\chapters@set{20261125T016500}{23}% hecke-operators-and-theta-lifts
\chapters@set{20261125T016700}{24}% divisibility-along-the-chain
\chapters@set{20261125T004000}{25}% denominators-and-sha
\chapters@set{20261125T004100}{26}% tate-twists-and-sha
\chapters@set{20261126T020430}{27}% the-minus-part-by-gauss-sums
\chapters@set{20261125T014100}{28}% congruence-modules
\chapters@set{20261125T014500}{29}% eisenstein-congruences-and-selmer-classes
\chapters@set{20261125T014700}{30}% pseudocharacters-and-the-length-bound
\chapters@set{20261125T011000}{31}% galois-representations-black-box
\chapters@set{20261126T030300}{32}% the-data
\chapters@set{20261126T020210}{33}% bessel-poincare-series-and-their-gram-operators
\chapters@set{20261126T030400}{34}% the-seesaw-identity
\chapters@set{20261126T030500}{35}% divisibility-of-the-composite-lift
\chapters@set{20261126T030600}{36}% theorem-a-short-forms
\chapters@set{20261126T030700}{37}% the-determinant-of-a-chain
\chapters@set{20261126T030800}{38}% the-cost-of-one-lift
\chapters@set{20261126T030900}{39}% denominators-of-projections-to-lifts
\chapters@set{20261126T031000}{40}% twists-with-few-congruences
\chapters@set{20261126T031100}{41}% theorem-c-and-the-main-inequality
\providecommand\chapnumber[1]{%
  \@ifundefined{chapters@#1}%
    {\PackageError{chapters}{Sheet #1 is not listed in the index}%
      {List it in 20261125T012000--grh-workbook__guide.org\MessageBreak
       and run scripts/make-chapters.py.}}%
    {\edef\thechap{\@nameuse{chapters@#1}}}}
\makeatother
\chapnumber{<sheet ID>}
\makeatletter
\providecommand\chapters@set[2]{\@namedef{chapters@#1}{#2}}
\chapters@set{20261125T001000}{1}% zeros-to-ramp-sums
\chapters@set{20261125T001100}{2}% inert-primes-and-estimate-b
\chapters@set{20261125T005000}{3}% galois-fourier-and-averaging
\chapters@set{20261125T005500}{4}% columns-from-a-sphere
\chapters@set{20261125T000000}{5}% slope-method
\chapters@set{20261126T030100}{6}% a-twist-for-each-prime
\chapters@set{20261125T003000}{7}% the-chain-and-the-square-root
\chapters@set{20261125T010000}{8}% rank-lemma-and-zero-estimate
\chapters@set{20261125T010500}{9}% poincare-series-and-near-orthogonality
\chapters@set{20261124T223000}{10}% elliptic-curves-for-the-proof
\chapters@set{20261125T009500}{11}% fock-model-and-theta-intertwiners
\chapters@set{20261124T223100}{12}% complex-multiplication-and-inert-reduction
\chapters@set{20261125T008000}{13}% one-cm-determinant
\chapters@set{20261124T214500}{14}% supersingular-theta-divisibility
\chapters@set{20261124T214600}{15}% the-isotropic-plane
\chapters@set{20261125T013700}{16}% pairings-on-a-formal-group-of-height-two
\chapters@set{20261125T013500}{17}% hecke-nondegeneracy-on-the-tree
\chapters@set{20261126T030200}{18}% the-bound-is-attained
\chapters@set{20261125T009000}{19}% kronecker-limit-and-chowla-selberg
\chapters@set{20261125T002000}{20}% definite-theta-lifts
\chapters@set{20261125T015500}{21}% siegel-weil-in-the-smallest-case
\chapters@set{20261125T002100}{22}% definite-theta-lifts-and-norms
\chapters@set{20261125T016500}{23}% hecke-operators-and-theta-lifts
\chapters@set{20261125T016700}{24}% divisibility-along-the-chain
\chapters@set{20261125T004000}{25}% denominators-and-sha
\chapters@set{20261125T004100}{26}% tate-twists-and-sha
\chapters@set{20261126T020430}{27}% the-minus-part-by-gauss-sums
\chapters@set{20261125T014100}{28}% congruence-modules
\chapters@set{20261125T014500}{29}% eisenstein-congruences-and-selmer-classes
\chapters@set{20261125T014700}{30}% pseudocharacters-and-the-length-bound
\chapters@set{20261125T011000}{31}% galois-representations-black-box
\chapters@set{20261126T030300}{32}% the-data
\chapters@set{20261126T020210}{33}% bessel-poincare-series-and-their-gram-operators
\chapters@set{20261126T030400}{34}% the-seesaw-identity
\chapters@set{20261126T030500}{35}% divisibility-of-the-composite-lift
\chapters@set{20261126T030600}{36}% theorem-a-short-forms
\chapters@set{20261126T030700}{37}% the-determinant-of-a-chain
\chapters@set{20261126T030800}{38}% the-cost-of-one-lift
\chapters@set{20261126T030900}{39}% denominators-of-projections-to-lifts
\chapters@set{20261126T031000}{40}% twists-with-few-congruences
\chapters@set{20261126T031100}{41}% theorem-c-and-the-main-inequality
\providecommand\chapnumber[1]{%
  \@ifundefined{chapters@#1}%
    {\PackageError{chapters}{Sheet #1 is not listed in the index}%
      {List it in 20261125T012000--grh-workbook__guide.org\MessageBreak
       and run scripts/make-chapters.py.}}%
    {\edef\thechap{\@nameuse{chapters@#1}}}}
\makeatother
\chapnumber{<sheet ID>}
\makeatletter
\providecommand\chapters@set[2]{\@namedef{chapters@#1}{#2}}
\chapters@set{20261125T001000}{1}% zeros-to-ramp-sums
\chapters@set{20261125T001100}{2}% inert-primes-and-estimate-b
\chapters@set{20261125T005000}{3}% galois-fourier-and-averaging
\chapters@set{20261125T005500}{4}% columns-from-a-sphere
\chapters@set{20261125T000000}{5}% slope-method
\chapters@set{20261126T030100}{6}% a-twist-for-each-prime
\chapters@set{20261125T003000}{7}% the-chain-and-the-square-root
\chapters@set{20261125T010000}{8}% rank-lemma-and-zero-estimate
\chapters@set{20261125T010500}{9}% poincare-series-and-near-orthogonality
\chapters@set{20261124T223000}{10}% elliptic-curves-for-the-proof
\chapters@set{20261125T009500}{11}% fock-model-and-theta-intertwiners
\chapters@set{20261124T223100}{12}% complex-multiplication-and-inert-reduction
\chapters@set{20261125T008000}{13}% one-cm-determinant
\chapters@set{20261124T214500}{14}% supersingular-theta-divisibility
\chapters@set{20261124T214600}{15}% the-isotropic-plane
\chapters@set{20261125T013700}{16}% pairings-on-a-formal-group-of-height-two
\chapters@set{20261125T013500}{17}% hecke-nondegeneracy-on-the-tree
\chapters@set{20261126T030200}{18}% the-bound-is-attained
\chapters@set{20261125T009000}{19}% kronecker-limit-and-chowla-selberg
\chapters@set{20261125T002000}{20}% definite-theta-lifts
\chapters@set{20261125T015500}{21}% siegel-weil-in-the-smallest-case
\chapters@set{20261125T002100}{22}% definite-theta-lifts-and-norms
\chapters@set{20261125T016500}{23}% hecke-operators-and-theta-lifts
\chapters@set{20261125T016700}{24}% divisibility-along-the-chain
\chapters@set{20261125T004000}{25}% denominators-and-sha
\chapters@set{20261125T004100}{26}% tate-twists-and-sha
\chapters@set{20261126T020430}{27}% the-minus-part-by-gauss-sums
\chapters@set{20261125T014100}{28}% congruence-modules
\chapters@set{20261125T014500}{29}% eisenstein-congruences-and-selmer-classes
\chapters@set{20261125T014700}{30}% pseudocharacters-and-the-length-bound
\chapters@set{20261125T011000}{31}% galois-representations-black-box
\chapters@set{20261126T030300}{32}% the-data
\chapters@set{20261126T020210}{33}% bessel-poincare-series-and-their-gram-operators
\chapters@set{20261126T030400}{34}% the-seesaw-identity
\chapters@set{20261126T030500}{35}% divisibility-of-the-composite-lift
\chapters@set{20261126T030600}{36}% theorem-a-short-forms
\chapters@set{20261126T030700}{37}% the-determinant-of-a-chain
\chapters@set{20261126T030800}{38}% the-cost-of-one-lift
\chapters@set{20261126T030900}{39}% denominators-of-projections-to-lifts
\chapters@set{20261126T031000}{40}% twists-with-few-congruences
\chapters@set{20261126T031100}{41}% theorem-c-and-the-main-inequality
\providecommand\chapnumber[1]{%
  \@ifundefined{chapters@#1}%
    {\PackageError{chapters}{Sheet #1 is not listed in the index}%
      {List it in 20261125T012000--grh-workbook__guide.org\MessageBreak
       and run scripts/make-chapters.py.}}%
    {\edef\thechap{\@nameuse{chapters@#1}}}}
\makeatother
\chapnumber{20261126T051600}
\renewcommand{\thesection}{\thechap.\arabic{section}}

\theoremstyle{plain}
\newtheorem{theorem}{Theorem}[section]
\theoremstyle{definition}
\newtheorem{problem}{Problem}[section]
\newtheorem{definition}[theorem]{Definition}
\theoremstyle{remark}
\newtheorem*{leanremark}{In the Lean formalization}

\newcommand{\Z}{{\mathbb{Z}}}
\newcommand{\Q}{{\mathbb{Q}}}
\newcommand{\C}{{\mathbb{C}}}
\newcommand{\F}{{\mathbb{F}}}
\newcommand{\Oc}{\mathcal O}
\newcommand{\ram}{\mathrm{ram}}
\newcommand{\Up}{\mathrm{Up}}
\newcommand{\Rt}{\mathrm{Rt}}
\newcommand{\Sch}{\mathrm{Sch}}
\newcommand{\dir}{\mathrm{dir}}
\newcommand{\col}{\mathrm{col}}
\newcommand{\ch}{\mathrm{ch}}
\newcommand{\Prim}{\mathrm{Prim}}
\newcommand{\dg}{\mathrm{diag}_N}
\newcommand{\elast}{e_{\mathrm{last}}}
\newcommand{\Dl}{D_{\mathrm{last}}}
\newcommand{\cdir}{c_{\mathrm{dir}}}
\newcommand{\cjet}{c_{\mathrm{jet}}}
\newcommand{\Tc}{\mathcal T}
\newcommand{\nb}[1]{\lVert #1 \rVert}
\newcommand{\tl}[1]{\texttt{#1}}
\newcommand{\us}{\_\allowbreak}
\setlength{\emergencystretch}{3em}

\title{The level prime and the ramified forms}
\author{}
\date{November 26, 2026}

\begin{document}
\maketitle
\refstepcounter{chap}\label{chap}

\paragraph{Prerequisites.}
Chapter~\ref{da-chap}: the scales and ranks (Definition~\ref{da-def:data}) and \(m\log b \le \log N \le 30m\log b\) (Problem~\ref{da-pr:MN}).  Chapter~\ref{ob-chap}: the CM type \(\Phi\), the source forms \(S\), their level, the operators \(T_w\), the ramified forms \(S_\ram\) and \(D_X = \dim_\C S\) (Section~\ref{ob-sec:source}), and the diagram (Definition~\ref{ob-def:diagram}).  The gauge \(\nb{\omega}_g\) of a Weyl vector and ``valuation at least'' (Definition~\ref{ob-def:gauge}).  Chapter~\ref{dl-chap}, Section~\ref{dl-sec:chain}: \(\Oc_\eta\) and \(S_\eta\), and Problem~\ref{dl-pr:gauge-chain}.  Chapter~\ref{ss-chap}: translations \(\mathfrak t \in \Tc\), the translated vacua \(v_j(\mathfrak t)\) and the definite lifts (Definition~\ref{ss-def:definite}); the pair gluing \(\mathrm{Gl}_i\), the glued tensor \(\mathrm{gl}(\mathfrak t)\), the restriction \(\mathrm{res}\) to \(Y_m\) and the composite lift \(B_{\vartheta,\Psi}\) (Definition~\ref{ss-def:composite}); chain data (Definition~\ref{ss-def:chain-datum}).  Chapter~\ref{d2-chap}, Section~\ref{d2-sec:kernel}: standard data of denominator \(\delta\) (Definition~\ref{d2-def:standard}), ``multiplies gauges by at most'' (Definition~\ref{d2-def:raises}), the jet scale \(\cjet\), statements (I), (II) and Problem~\ref{d2-pr:scaling}.  The colength of Chapter~\ref{sd-chap} (before Theorem~\ref{sd-thm:colength}), with the same definition over \(\Oc_\eta\).  Standard fillings and the Weyl modules of \(\mathrm{GL}_k\) \cite[\S8]{Fu}.

\medskip
At the level prime the gain is \(0\), but the glued pairs of the composite lift carry poles at \(N\): their data translate the vacuum by \(-y/N\).  How much \(N\)-adic valuation does the composite lift lose there, and on how large a module of forms?  This is the case \(p = N\) of Theorem~\ref{dl-thm:one-prime}; Section~\ref{dl-sec:sum} uses it in the form of Theorem~\ref{thm:level} below.  On ramified forms every glued pair becomes two direct lifts after an operator \(\Up\), and \(\Up^{m/2}\) is integral, after a change of base point, on a submodule of small colength.

\section{Ramified forms}\label{sec:ramified}

Recall from Section~\ref{ob-sec:source} that for each prime \(w \mid \pi_N\) of \(E\) right translation by the idele \(z_w\) is an operator \(T_w\) on \(S\) with \(T_w^{N-1} = 1\), and \(S_\ram = \bigcap_{w \mid \pi_N}\ker\sum_{i=0}^{N-2}T_w^i\).

\begin{problem}\label{pr:ram-image}
Let \(F_1, \dots, F_l\) be commuting operators on a finite-dimensional complex space with \(F_i^n = 1\).  Then \(\bigcap_i \ker\sum_{s<n}F_i^s = \operatorname{im}\prod_i(1 - F_i)\).  Hence \(S_\ram = \operatorname{im}\prod_{w \mid \pi_N}(1 - T_w)\), as in Problem~\ref{ss-pr:ram-image}.  For the cyclic shift \(F\) of \(\C^4\), \(F e_s = e_{s+1}\), the vector \((1-F)e_0 = (1,-1,0,0)\) is killed by \(\sum_{s<4}F^s\).
\end{problem}

\begin{problem}\label{pr:sign}
Let \(N = 3\) and let \(F\) act on \(\C\) through the sign character of \(\F_3^\times\) at its generator.  Then \(1 + F = 0\), \(1 - F = 2\), so the line is ramified and \(\lvert 1 - F\rvert^2 = 4\).
\end{problem}

\section{The places above \texorpdfstring{\(N\)}{N} and the direct lifts}\label{sec:places}

Fix a ring map \(\eta\colon \C \to \C_N\).  For \(\tau \in \Phi\) let \(w_\tau\) be the place of \(E\) above \(N\) at which \(\eta \circ \Phi_\tau\) is continuous (the prime \(v_\tau\) of Section~\ref{dl-sec:exceptional}, for \(p = N\); we write \(w\) for places above \(N\), as in Section~\ref{ob-sec:source}).  Let \(\Sigma_N\) (a notation of this chapter) be the set of the conjugates \(\bar w_\tau\).  As \(N \equiv 1 \pmod{q_2^2}\), \(N\) splits completely in \(E\); so \(\tau \mapsto \bar w_\tau\) is a bijection from \(\Phi\) onto \(\Sigma_N\), which holds exactly one place of each conjugate pair above \(N\); each residue field is \(\F_N\).  A \emph{residue vector} on \(X_j\) is a family \(y = (y_w)_{w \in \Sigma_N}\) of vectors \(y_w \in \F_N^{n_j}\) in the basis of \(X_j\); it is \emph{primitive} if every \(y_w \ne 0\); \(\elast\) is the last basis vector at every place.

\begin{definition}\label{def:base-point}
Let \(\dg\) be the matrix \(\mathrm{diag}(1, \dots, 1, N)\) at every place of \(\Sigma_N\).  The \emph{base point} \(\delta_j\) is the finite adelic point of the unitary group of \(X_j\), supported above \(N\), whose component at \(w \in \Sigma_N\) is \(\mathrm{diag}(1, \dots, 1, N)\) if \(w \mid \pi_N\) and \(1\) otherwise.  Put \(\delta_j^\vee := \delta_j\,\delta_{j,\Sigma}^{-1}\), where \(\delta_{j,\Sigma}\) is the point with component \(\dg\) at \(\Sigma_N\).
\end{definition}

For a matrix \(g_1\) over \(\prod_{w \in \Sigma_N}E_w\) let \(\Rt_j(g_1)\) be right translation of forms on \(X_j\) by \(\delta_j g_1\,\delta_j^{-1}\): matrices at \(\Sigma_N\) are read through \(\delta_j\).  For a function \(f\) of \(n_{j+1} \times n_j\) matrices at \(\Sigma_N\) and a finite datum \(v\) of \(\mathrm{Hom}(X_j, X_{j+1})\), let \(\Sch(f)v\) put \(f\) on \(v\) at \(\Sigma_N\) in the Schr\"odinger model, and let \(\Theta^X_j[f, v]\) (in this chapter) be the definite lift \(X_j \to X_{j+1}\) for the twist \(\vartheta\) (lift \(j+1\) of Definition~\ref{ss-def:definite}; here, as in Definition~\ref{dl-def:restricted}, a lift carries the index of its source space) whose datum is \(\Sch(f)v\) translated by \((\delta_j, \delta_{j+1})\) under the Weil representation.  A function of residue matrices is read on the integral matrices through reduction modulo \(N\) and extended by \(0\).

\begin{definition}\label{def:direct}
The \emph{direct datum} \(\dir\) is the function \(\col_{\elast}\), where \(\col_y(T) = 1\) if \(T\) is integral and its last column reduces to \(y\), and \(\col_y(T) = 0\) otherwise.  The \emph{direct lift} \(j\) with datum \(v\) is \(\Theta^X_j[\dir, v]\).
\end{definition}

\begin{definition}\label{def:up}
Let \(\psi_N(c) := \psi_0(-c/N)\) on \(\prod_{w \in \Sigma_N}\F_N\), with \(\psi_0\) the standard character of the finite adeles.  For a residue vector \(y\) on \(X_0\) whose last entry is a unit at every place, let \(\mathfrak g_y\) be an integral invertible matrix whose last column and last row reduce to \(y\).  On forms on \(X_0\),
\[
  \Up := \Rt_0(\dg^{-1}) \sum_{y} \psi_N(y_{\mathrm{last}})\,\Rt_0(\mathfrak g_y),
\]
the sum over these \(y\).
\end{definition}

A form \(\varphi\) on \(X_j\) is \emph{ramified at \(\Sigma_N\)} (in this chapter) if for each \(w \in \Sigma_N\) there are a form \(\psi\) and a unit scalar \(c\) at \(\Sigma_N\), integral, with \(c \equiv 1 \bmod N\) at the places other than \(w\), such that \(\varphi = \psi - \Rt_j(c)\psi\).

\begin{problem}\label{pr:ram-vector}
Every \(\varphi \in S_\ram\) is ramified at \(\Sigma_N\), with \(\psi \in S\).
\end{problem}

Read from the proof: the following \emph{lift law} for \(\Theta^X_j[\,\cdot\,, v]\), for every datum \(v\).  (i)~For integral invertible matrices \(g_1\) of size \(n_j\) and \(g_2\) of size \(n_{j+1}\) at \(\Sigma_N\), with reductions \(\bar g_1\), \(\bar g_2\), and every function \(f\) of residue matrices, \(\Rt_{j+1}(g_2)\circ\Theta^X_j[T \mapsto f(\bar g_2T\bar g_1), v]\circ\Rt_j(g_1) = \Theta^X_j[f, v]\).  (ii)~\(\Rt_{j+1}(\dg^{-1})\circ\Theta^X_j[\mathrm{row}, v] = \varepsilon'_j\,\Theta^X_j[\dir, v]\circ\Rt_j(\dg^{-1})\), where \(\mathrm{row}\) is the indicator of the integral matrices whose last row reduces to \(\elast\), and \(\varepsilon'_j\) (in this chapter) is the value at \(\dg\) on \(X_j\) of the character \(g \mapsto \chi(\det g)\lvert\det g\rvert^{-n_{j+1}/2}\), divided by the value at \(\dg\) on \(X_{j+1}\) of \(g \mapsto \chi(\det g)\lvert\det g\rvert^{-n_j/2}\), with \(\chi\) the base change to \(E\) of the character of the lift \(X_j \to X_{j+1}\).  Here and below \(\lvert\det g\rvert := \prod_{w \in \Sigma_N}\lvert\det g_w\rvert_w\), with \(\lvert N\rvert_w = N^{-1}\), and \(\chi(\det g)\) reads \(\det g\) as an idele.

\begin{problem}\label{pr:non-primitive}
Let \(w \in \Sigma_N\) and let \(f\) be a function of residue matrices with \(f(cT) = f(T)\) for every unit \(c \in \prod_{\Sigma_N}\F_N^\times\) equal to \(1\) off \(w\).  Then \(\Theta^X_j[f, v]\varphi = 0\) for every \(\varphi\) ramified at \(\Sigma_N\).  In particular \(\Theta^X_j[\col_y, v]\varphi = 0\) if \(y\) vanishes at some place.
\end{problem}

Let \(\ch_y(T) := \psi_N\bigl((\text{last row of } T)\cdot y\bigr)\) on integral matrices, \(0\) elsewhere.

\begin{problem}\label{pr:primitive}
\(\mathrm{GL}_n(\F_N)\) is transitive on \(\F_N^n \smallsetminus 0\), a set of \(N^n - 1\) elements; so there are \(\#\Prim_n = (N^n-1)^f\) primitive residue vectors in rank \(n\), a number prime to \(N\).  For \(\varphi\) ramified at \(\Sigma_N\) and data \(v, v'\),
\[
  \sum_y \Theta^X_{j+1}[\ch_y, v']\,\Theta^X_j[\col_y, v]\,\varphi = \#\Prim_{n_{j+1}}\cdot\Theta^X_{j+1}[\ch_{\elast}, v']\,\Theta^X_j[\dir, v]\,\varphi ,
\]
the sum over all residue vectors on \(X_{j+1}\).  For \(N = 3\), \(n = 2\), \(f = 1\): eight vectors, one orbit.
\end{problem}

\emph{All lifts become direct.}  Two further statements are read from the proof.  First, on data fixed by the projection on the vacuum above \(N\), the double lift \(X_{2i} \to X_{2i+2}\) with the glued datum \(\mathrm{Gl}_i(v \otimes v')\) is \(\varepsilon_i\,\Rt_{2i+2}(\dg^{-1})\sum_y\Theta^X_{2i+1}[\ch_y, v']\,\Theta^X_{2i}[\col_y, v]\), where \(\varepsilon_i\) (in this chapter) is the value of the character of the lift \(X_{2i+1} \to X_{2i+2}\), formed with \(\lvert\det\rvert^{-n_{2i+1}/2}\), at \(\mathrm{diag}(1,\dots,1,N)\) on \(X_{2i+2}\) at the places of \(\Sigma_N\) prime to \(\pi_N\), times the value of the character of \(X_{2i} \to X_{2i+1}\), also formed with \(\lvert\det\rvert^{-n_{2i+1}/2}\), at the same matrix on \(X_{2i}\) at the places dividing \(\pi_N\).  Second, the following, whose proof sorts the terms of \(\Theta^X_{2i+1}[\ch_{\elast}, v']\) by the last row and uses the lift law:
\emph{there is \(\cdir \in \C\) (in this chapter) such that for all data \(v_j\) and all \(\varphi \in S_\ram\), the chain of the \(m/2\) double lifts with the glued data of the projections of the \(v_j\) on the vacuum above \(N\) sends \(\varphi\) to \(\cdir\) times the chain of the \(m\) direct lifts with data \(v_j\) applied to \(\Up^{m/2}\varphi\).  Here}
\[
  \cdir = \prod_{i < m/2}\Bigl(\varepsilon_i\,\#\Prim_{n_{2i+1}}\prod_{l < 2i+2}\varepsilon'_l\Bigr) .
\]
\emph{Status}: read from the proof.  Problems~\ref{pr:non-primitive} and~\ref{pr:primitive} are its first steps.  We also read: for \(g\) equal to \(\mathrm{diag}(1, \dots, 1, N)\) at a set \(\Sigma'\subseteq\Sigma_N\) and \(1\) at the other places of \(\Sigma_N\), the value at \(g\) of the character of any of the lifts, formed with \(\lvert\det\rvert^{-r/2}\) for any \(r \ge 0\), has \(N\)-adic absolute value at most \(N^{2b\,\#\Sigma'}\) after \(\eta\).

\begin{problem}\label{pr:scalar}
\(\lvert\eta(\varepsilon_i)\rvert \le N^{2fb}\), \(\lvert\eta(\varepsilon'_l)\rvert \le 1\) and \(\lvert\eta(\#\Prim_n)\rvert \le 1\); hence \(\lvert\eta(\cdir)\rvert \le N^{mfb}\).
\end{problem}

\section{The colength of the forms integral after \texorpdfstring{\(\Up^{m/2}\)}{Up}}\label{sec:colength}

A form \(\varphi\) is \emph{integral at} a point \(h\) if \(\nb{\varphi(g)}_{gh} \le 1\) for every \(g\); so \(S_\eta\) is the module of source forms integral at \(1\).  Write \(S_\eta(h)\) (in this chapter) for the \(\Oc_\eta\)-module of source forms integral at \(h\); so \(S_\eta = S_\eta(1)\).  By Section~\ref{ob-sec:source} the values of source forms lie in \(\bigotimes_{\tau \in \Phi_a}\mathrm{Sym}^a\C^k\), whose basis is the families \(I = (z^{\alpha_\tau})_{\tau \in \Phi_a}\) of monomials of degree \(a\) in \(k\) variables; let \(e_I := \sum_\tau \alpha_{\tau,k}\) be the degree of \(I\) in the last variable.

\begin{problem}\label{pr:sym2}
Let \(k = a = 2\), one place, \(N = 3\).  The monomials \(z_1^2, z_1z_2, z_2^2\) have degrees \(0, 1, 2\) in \(z_2\); the change of lattice multiplies them by \(\mathrm{diag}(1, 3, 9)\), of determinant \(3^3\), and the colength of the span of \(3^{\alpha_2}z^\alpha\) in the span of the \(z^\alpha\) over \(\Z_3\) divides \(3^3\).  The mean degree in \(z_2\) is \(1 = a/k\).
\end{problem}

A source form is determined by its values at representatives of the double cosets of its level, and these values are arbitrary, since (read from the proof) the stabilizers of the double cosets are trivial.  So the pairs \((c, I)\) of a double coset and a basis vector are \(D_X\) in number.  Let \(\Dl := \sum_{(c, I)} e_I\).

\begin{problem}\label{pr:symmetry}
The degree in the \(u\)-th variable, summed over the monomials of degree \(a\) in \(k\) variables, does not depend on \(u\), and \(k\) times it is \(a\) times the number of these monomials.
\end{problem}

\begin{problem}\label{pr:degree}
\(k\,\Dl = \#\Phi_a\,a\,D_X \le fa\,D_X\).
\end{problem}

\begin{definition}\label{def:SN}
\(S_N\) is the \(\Oc_\eta\)-module of the forms \(\varphi\) on \(X_0\) such that \(\Up^{m/2}\varphi\) is a source form integral at \(\delta_0\).
\end{definition}

Read from the proof: every double coset has a representative \(r_c\) such that \(r_c\delta_0^\vee\) and \(\delta_{0,\Sigma}^{-1}r_c\delta_0\) are integral above \(N\), that is, lie in \(\mathrm{GL}_{k}(\Oc_{E_w})\) at every place \(w\) above \(N\).

\begin{problem}\label{pr:colength-model}
\(S_\eta \subseteq S_\eta(\delta_0^\vee)\).  The coordinates of \(\varphi(r_c)\) are integral for \(\varphi \in S_\eta(\delta_0^\vee)\), and a source form whose coordinate \((c, I)\) at \(r_c\) is divisible by \(N^{e_I}\) for all \((c, I)\) lies in \(S_\eta(\delta_0)\).  Hence the colength of \(S_\eta(\delta_0)\) in \(S_\eta(\delta_0^\vee)\) divides \(N^{\Dl}\).
\end{problem}

\begin{problem}\label{pr:up-integral}
\((\Up\varphi)(g) = \sum_y\psi_N(y_{\mathrm{last}})\,\varphi(g\,\delta_0^\vee\mathfrak g_y\delta_0^{-1})\); hence \(\Up\) maps \(S_\eta(\delta_0)\) into forms integral at \(\delta_0^\vee\).  As \(\Up\) preserves \(S\) (read from the proof), the colength in \(S_\eta\) of the forms \(\varphi\) with \(\Up^i\varphi \in S_\eta(\delta_0)\) divides \(N^{(i+1)\Dl}\), for every \(i \ge 0\).
\end{problem}

\begin{problem}\label{pr:exponent}
\((m/2 + 1)\,a \le m b k\), and so \((m/2 + 1)\Dl \le mfb\,D_X\).
\end{problem}

Problems~\ref{pr:degree}, \ref{pr:up-integral} with \(i = m/2\), and~\ref{pr:exponent} prove part~(a) of Theorem~\ref{thm:level} below, with \(e = (m/2+1)\Dl\).

\section{Direct lifts lose little}\label{sec:direct}

The Fourier expansion of the direct datum has poles at \(N\), which an element of \(K\) must clear.

\begin{problem}\label{pr:sqrt-7}
Let \(K = \Q(\sqrt{-7})\), \(\theta_K = (1+\sqrt{-7})/2\), \(N = 11\).  Then \(z^2 - z + 2\) has the roots \(5, 7\) modulo \(11\); with \(c_0 = 7\), \((1 - 2c_0)^2 \equiv -7\) and \((1 - 2c_0)\cdot 7 \equiv 8 \pmod{11}\).  So \(\delta_K := \theta_K - c_0\) lies in the prime \(\mathfrak p = (11, \delta_K)\), and \(\bar{\delta}_K\cdot 7\) is a unit at \(\mathfrak p\).
\end{problem}

\begin{problem}\label{pr:delta}
Let \(\theta_K = (1 + \sqrt{-d})/2\), let \(c_0 \in \Z\) represent \(\theta_K\) modulo the prime \(\mathfrak p\) of \(K\) below \(w_\tau\) (the same for all \(\tau\)), and \(\delta_K := \theta_K - c_0\).  Then \(\delta_K \in w_\tau\) for every \(\tau\), and \(\lvert\eta(\bar{\delta}_K d)\rvert = 1\).
\end{problem}

Read from the proof: \emph{for \(j < m\), every translation \(\mathfrak t\) and every \(\delta_K \in \Oc_K\) lying in all the \(w_\tau\), the datum \(N^{2fn_jn_{j+1}}\,\Sch(\dir)v_j(\mathfrak t)\) is standard of denominator \(\delta_K\)}; the translates of the vacuum in the expansion of \(\dir\) have simple poles at the \(w_\tau\).  Recall from Section~\ref{d2-sec:kernel} that the kernel of the lift \(X_j \to X_{j+1}\) is the theta kernel at scale \(\pi_M\cjet\) and \(\lvert\eta\pi_M\rvert \le 1\).  A datum \emph{multiplies gauges from \(\delta_j\) to \(\delta_{j+1}\) by at most \(c\)} if the inequality of Definition~\ref{d2-def:raises} holds with \(\nb{\cdot}_{g'\delta_{j+1}}\) and \(\nb{\cdot}_{g\delta_j}\) in place of \(\nb{\cdot}_{g'}\) and \(\nb{\cdot}_g\).

\begin{problem}\label{pr:direct-step}
Let \(0 \ne \delta_K \in \Oc_K\) with \(\lvert\eta\pi_M\rvert \le \lvert\eta(\bar{\delta}_K d)\rvert\).  Every standard datum of denominator \(\delta_K\) multiplies the gauges of the kernel of the lift \(X_j \to X_{j+1}\) by at most \(1\).  Hence, for \(j < m\), the datum of the direct lift \(j\) with \(v_j(\mathfrak t)\) multiplies gauges from \(\delta_j\) to \(\delta_{j+1}\) by at most \(N^{2fn_jn_{j+1}} \le N^{fb}\).
\end{problem}

Three more statements are read from the proof.  For \(\mathfrak t\) in the support of a chain datum, \(v_j(\mathfrak t)\) is fixed by the projection on the vacuum above \(N\).  The direct datum commutes with the action of the levels of \(X_j\) and \(X_{j+1}\); so for a chain datum \(\Psi\) for \(\vartheta\), if every \(\Sch(\dir)v_j(\mathfrak t)\), translated by \((\delta_j, \delta_{j+1})\), multiplies gauges from \(\delta_j\) to \(\delta_{j+1}\) by at most \(N^{r_j}\), then \(\sum_{\mathfrak t}\Psi(\mathfrak t)\) times the chain of the \(m\) direct lifts with data \(v_j(\mathfrak t)\) maps source forms integral at \(\delta_0\) to forms with \(\nb{\cdot(g)}_{g\delta_m} \le N^{r_0 + \dots + r_{m-1}}\) for every \(g\), as in Problem~\ref{dl-pr:gauge-chain}.  And if a form on \(X_m\) satisfies \(\nb{\varphi(g)}_{g\delta_m} \le C\) for every \(g\), then its restriction \(\mathrm{res}\,\varphi\) to \(Y_m\) has valuation at least \(-\log_N C\): the branching map is integral, and its scalar has the absolute value of the determinant factor of \(\delta_m\).  The branching rule is classical \cite[Thm.~8.1.1]{GW}; the integrality of the branching map and the rest are read from the proof.

\begin{problem}\label{pr:two-costs}
Let \(\Psi\) be a chain datum for a twist \(\vartheta\) and \(\varphi \in S_\ram \cap S_N\).  Then \(B_{\vartheta,\Psi}\varphi = \cdir\,\mathrm{res}\bigl(\sum_{\mathfrak t}\Psi(\mathfrak t)\,(\text{chain of direct lifts with } v_j(\mathfrak t))(\Up^{m/2}\varphi)\bigr)\), and it has valuation at least \(-2mfb\).  Moreover \(2mfb\log N \le 4fa(\log b)^2\log x_1\).
\end{problem}

Problems~\ref{pr:scalar}, \ref{pr:direct-step} and~\ref{pr:two-costs} prove part~(b).

\begin{theorem}\label{thm:level}
Let \(\eta\colon\C\to\C_N\) be a ring map.
\begin{enumerate}
\item[(a)] There is an integer \(0 \le e \le mfb\,D_X\) such that the colength of \(S_N\) in \(S_\eta\) divides \(N^e\).
\item[(b)] For every twist \(\vartheta\), every chain datum \(\Psi\) for \(\vartheta\) and every \(\varphi \in S_\ram \cap S_N\), the form \(B_{\vartheta,\Psi}\varphi\) has valuation at least \(-2mfb\).
\end{enumerate}
\end{theorem}

\paragraph{What has been established.}  Theorem~\ref{thm:level}, the case \(p = N\) of Theorem~\ref{dl-thm:one-prime}, with \(S' = S_N \cap S_\eta\) (the combinations in (a) of elements of \(S_\eta\) stay in \(S_\eta\)), \(l = e\) and \(\mathrm{loss}(N) = 2mfb\); Problem~\ref{pr:two-costs} bounds its term in the sum of the errors.  Taken as given: the statements read from the proof, listed in the remark, and the branching rule.

\begin{leanremark}
\begin{sloppypar}
In the Lean proof \cite{GL} (namespace \tl{New.} omitted), Theorem~\ref{thm:level}(a) is \tl{Data.exists\us colengthDvd\us upIntegralForms\us pow} and (b) is \tl{Data.gaugeLe\us compositeLift\us at\us N}.  Definitions~\ref{def:base-point}, \ref{def:direct}, \ref{def:up}, \ref{def:SN} are \tl{Data.diagPointN} (with \tl{Data.dualPointN}), \tl{IsSplitLevel.directDatum}, \tl{IsSplitLevel.upOperator} and \tl{Data.upIntegralForms}; Problem~\ref{pr:ram-image} is \tl{Data.ramifiedSourceForms\us eq\us range}, Problem~\ref{pr:colength-model} is \tl{Data.colengthDvd\us integralFormsAt\us diagPointN}, and the statement that all lifts become direct is \tl{Data.doubleLiftChain\us gluedDoubleLiftDatum\us eq\us smul\us chain} (one pair: \tl{IsSplitLevel.LiftLaw.glued\us eq\us direct\us direct\us upOperator}); the integral branching map is \tl{exists\us branching\us matrix}.  The other statements read in the text are: the lift law, \tl{Data.liftLaw\us liftAtN} (\(\varepsilon'_j\) is \tl{Data.upCommutationFactor}); the gluing at \(N\), \tl{Data.liftInPairs\us gluedDoubleLiftDatum} (\(\varepsilon_i\) is \tl{Data.gluedFactor}); the bound on the character values, \tl{Data.norm\us detCharAt\us diagN\us le}; the trivial stabilizers, \tl{Data.classAut\us levelX\us zero\us eq\us bot}; the representatives \(r_c\), \tl{Data.exists\us mul\us dualPointN\us mem\us integralAbove}; that \(\Up\) preserves \(S\), \tl{Data.upOperator\us mem\us sourceForm}; the Fourier step, a step of the proof of \tl{Data.schrodingerOpAt\us directDatum\us mem\us raises}; the translated vacua at \(N\), \tl{Data.vacuumProjFinPart\us translateDatum}; the invariance and the bound along the chain, \tl{Data.omega\us schrodingerOpAtN\us directDatum} and \tl{Data.IsWeakChainDatum.gaugeLe\us chain}; the restriction at \(\delta_m\), \tl{Data.GaugeLe.restrictToY\us of\us diag}.
\end{sloppypar}
\end{leanremark}

\section*{Hints}

\subsection*{\S\ref{sec:ramified}}
\begin{description}
\item[\ref{pr:ram-image}] As in Problems~\ref{ss-pr:mean} and~\ref{ss-pr:mean-several}: \(1 - F^n = (1 - F)\sum_{s<n}F^s\); decompose into joint eigenspaces.
\end{description}

\subsection*{\S\ref{sec:places}}
\begin{description}
\item[\ref{pr:ram-vector}] Problem~\ref{pr:ram-image} with one factor.  For \(w \in \Sigma_N\) let \(u\) be the one of \(w, \bar w\) that divides \(\pi_N\); \(z_u\) is \(1\) off \(u, \bar u\), and a scalar commutes with \(\delta_0\).
\item[\ref{pr:non-primitive}] Lift law (i) with \(g_1\) the scalar \(c\) and \(g_2 = 1\).
\item[\ref{pr:primitive}] Write \(y = \bar g\,\elast\) with \(g\) integral invertible, and apply the lift law (i) to both lifts; Problem~\ref{pr:non-primitive} for the other \(y\).
\item[\ref{pr:scalar}] Each \(\tau\) has \(\bar w_\tau\) in exactly one of the two sets in \(\varepsilon_i\).  For \(\varepsilon'_l\): \(\chi(\det\dg)\) is the same on both spaces, and \(\lvert\eta(\sqrt N)\rvert = N^{-1/2}\); \(n_l \le n_{l+1}\).
\end{description}

\subsection*{\S\ref{sec:colength}}
\begin{description}
\item[\ref{pr:sym2}] \(\mathrm{Sym}^2\) of \(\mathrm{diag}(1, 3)\).  Over a valuation ring one condition ``this coordinate is divisible by \(z\)'' costs a colength dividing \(z\): use a member of the family whose coordinate divides the others'.
\item[\ref{pr:symmetry}] Exchanging two variables permutes the monomials; sum over \(u\).
\item[\ref{pr:degree}] Sum Problem~\ref{pr:symmetry} over the places of \(\Phi_a\), one at a time.
\item[\ref{pr:colength-model}] At \(\tau\), \(\delta_0^\vee\) acts through \(\mathrm{diag}(1, \dots, 1, N^{-1})\) or \(1\), and \(\delta_{0,\Sigma}\) through \(\mathrm{diag}(1, \dots, 1, N)\); the last row of an element of the level is \(\equiv (0, \dots, 0, 1) \bmod \pi_N\), so \(\delta_0^{-1}k\delta_0\) is integral above \(N\).  Every \(g\) is \(\gamma r_ck\) with \(\gamma\) rational and \(k\) in the level, and \(\nb{\gamma\omega}_{\gamma h} = \nb{\omega}_h\).  For the colength, impose the divisibility of the \((c, I)\) coordinates one at a time, as in the hint to Problem~\ref{pr:sym2}.
\item[\ref{pr:up-integral}] \(\mathfrak g_y\) is integral above \(N\), so \(\nb{\omega}_{h\mathfrak g_y} = \nb{\omega}_h\); and \(\lvert\psi_N\rvert = 1\).  Induction on \(i\); colengths multiply along a chain of modules.
\item[\ref{pr:exponent}] \(a = xb\), \(\sqrt x \le m\), \(k = r + 1 - m \ge x^5 - 2\sqrt x\); Problem~\ref{pr:degree}.
\end{description}

\subsection*{\S\ref{sec:direct}}
\begin{description}
\item[\ref{pr:sqrt-7}] \(\theta_K^2 - \theta_K + 2 = 0\); \(\bar{\delta}_K = 1 - \theta_K - c_0\).
\item[\ref{pr:delta}] \(\Phi_\tau\) restricts to one embedding of \(K\); \((1 - 2\theta_K)^2 = -d\), and \(N \nmid d\).
\item[\ref{pr:direct-step}] Statements (I), (II) of Section~\ref{d2-sec:kernel} and Problem~\ref{d2-pr:scaling}, with Problem~\ref{pr:delta}; the entries of \(\mathrm{schur}_\beta(z)\) are homogeneous.  Then \(\lvert\eta(N)\rvert = N^{-1}\), \(n_j \le n_m\), and the scale inequalities give \(2n_m^2 \le b\).
\item[\ref{pr:two-costs}] The chain of double lifts is linear in the glued tensor; Problems~\ref{pr:scalar} and~\ref{pr:direct-step}; the values \(\Psi(\mathfrak t)\) are integral at \(N\).  For the last claim use Problem~\ref{da-pr:MN}, \(m \le 2\sqrt x\) and \((\log b)\log x_1 \ge 60\).
\end{description}

\begin{thebibliography}{GW}
\bibitem[Fu]{Fu} W.~Fulton, \emph{Young Tableaux}, London Math. Soc. Student Texts 35, Cambridge University Press, 1997.
\bibitem[GL]{GL} L.~Sernau, \emph{GRH for Dirichlet \(L\)-functions, in Lean~4}, a shorter rewriting of the Lean proof, text by Claude, \url{https://github.com/lsernau/grh-lean}, commit \texttt{72fead1}, readable edition \texttt{grh/}.
\bibitem[GW]{GW} R.~Goodman and N.~R.~Wallach, \emph{Symmetry, Representations, and Invariants}, Graduate Texts in Mathematics 255, Springer, 2009.
\end{thebibliography}

\end{document}
